% migrate-to: execution-market-making, market-making-lob
\subsection{Inventory Models}

\begin{principle} \hlt{Inventory Spread Component}\\
Inventory risk is compensation for holding unwanted positions while waiting for offsetting flow.\\
It is separate from adverse selection: inventory risk can exist even when no trader is better informed, because prices can move before the dealer unwinds.
\end{principle}

\begin{taxonomy} \hlt{Inventory Cost Segmentation}
\begin{enumerate}[label=\roman*.]
\setlength{\itemsep}{0pt}
\item \hlt{price risk}: mark-to-market loss while inventory is held;
\item \hlt{funding risk}: financing cost of long inventory and borrow or recall risk for shorts;
\item \hlt{capital risk}: balance-sheet, margin, and risk-limit consumption;
\item \hlt{hedging cost}: basis risk and transaction costs from reducing exposure elsewhere.
\end{enumerate}
\end{taxonomy}

% assign: market-making-lob
\begin{model} \hlt{Inventory-Skew Quoting}\\
Let $q_t$ be inventory, $q^\star$ the target inventory, and $r_t=m_t-\kappa(q_t-q^\star)$ the dealer's reservation price.\\
Quotes can be written as $b_t=r_t-\frac{s_t}{2}$, $a_t=r_t+\frac{s_t}{2}$.\\
When $q_t>q^\star$, the dealer lowers the reservation price to attract buyers and discourage sellers; when $q_t<q^\star$, the dealer raises it to attract sellers and discourage buyers.
\end{model}

\begin{caveat} \hlt{Inventory vs Information Identification}\\
Inventory control and adverse selection can produce the same quote movement after a buy or sell.\\
Identification needs additional information: dealer inventory, hedges, client type, future midpoint revision, or whether quotes revert after inventory is restored.
\end{caveat}
